% Complete Section 3.1, original PDF 31-32 / printed 21-22.
% block: unbounded-definition-heading page:31 type:section
\subsection{\BT{Unbounded operators}{无界算子}}\label{sec:unbounded-operators}

% block: unbounded-domain-definition pages:31,32 type:ordinary-paragraph
\BT{When considering unbounded operators, one must allow for operators whose domain may not be the entire Hilbert space but only a subspace. Thus, by an \emph{unbounded operator} $T$ from a Hilbert space $X$ to a Hilbert space $Y$, we mean a linear operator mapping a subspace $D(T)$ of $X$, called the domain of $T$, into $Y$. The operator $T$ is not necessarily bounded and the domain $D(T)$ is not necessarily a closed subspace. Such operators are called unbounded operators on Hilbert space, even though the case of a bounded linear operator defined on all of $X$ is included. The more accurate term \emph{not-necessarily-everywhere-defined, not-necessarily-bounded linear operators} would be a bit unwieldy! We shall almost always be concerned with \emph{densely defined} operators, meaning that the domain $D(T)$ is a dense subspace of $X$. A typical example to keep in mind is $X=L^2(\Omega)$ for some domain $\Omega\subset\mathbb R^n$, $Y=L^2(\Omega;\mathbb R^n)$, the domain $D(T)$ is the Sobolev space $H^1(\Omega)$, which is a dense subspace of $L^2(\Omega)$, and the operator $T$ is the gradient. That is, $T\phi=\operatorname{grad}\phi\in L^2(\Omega;\mathbb R^n)$ for all $\phi\in H^1(\Omega)$. A variation is the operator $\widetilde T$ with $D(\widetilde T)=\mathring H^1(\Omega)$, which is also dense in $L^2$, again with $\widetilde T\phi=\operatorname{grad}\phi$, but with $\widetilde T\phi$ only defined for $\phi\in\mathring H^1$. Other possibilities would be to define a gradient operator whose domain consists of $C^1$ functions on $\overline\Omega$, or $C^\infty$ functions, or $C^\infty$ functions with compact support inside $\Omega$. Although we might refer to any one of these operators as the gradient, they are different, and it is important to always be clear about the domain. If we say that two unbounded operators $S$ and $T$ are equal, we mean both that $D(S)=D(T)$ and that the actions of the operators coincide.}{考虑无界算子时, 必须允许算子的定义域不是整个 Hilbert 空间, 而只是一个子空间. 因此, 所谓从 Hilbert 空间 $X$ 到 Hilbert 空间 $Y$ 的\emph{无界算子} $T$, 是指将 $X$ 的一个子空间 $D(T)$ 映入 $Y$ 的线性算子, 其中该子空间称为 $T$ 的定义域. 算子 $T$ 不一定有界, 定义域 $D(T)$ 也不一定是闭子空间. 这类算子称为 Hilbert 空间上的无界算子, 尽管其中也包含在整个 $X$ 上定义的有界线性算子. 更准确的名称\emph{不一定处处有定义且不一定有界的线性算子}就有些冗长了! 我们几乎总是关注\emph{稠密定义}的算子, 即定义域 $D(T)$ 是 $X$ 的稠密子空间. 一个值得牢记的典型例子是: 对某个区域 $\Omega\subset\mathbb R^n$, 取 $X=L^2(\Omega)$, $Y=L^2(\Omega;\mathbb R^n)$, 定义域 $D(T)$ 为 Sobolev 空间 $H^1(\Omega)$, 它是 $L^2(\Omega)$ 的稠密子空间, 而算子 $T$ 为梯度. 也就是说, 对所有 $\phi\in H^1(\Omega)$, 有 $T\phi=\operatorname{grad}\phi\in L^2(\Omega;\mathbb R^n)$. 一个变体是算子 $\widetilde T$, 其定义域 $D(\widetilde T)=\mathring H^1(\Omega)$ 同样在 $L^2$ 中稠密, 且仍有 $\widetilde T\phi=\operatorname{grad}\phi$, 但 $\widetilde T\phi$ 只对 $\phi\in\mathring H^1$ 定义. 还可以定义其他梯度算子, 其定义域由 $\overline\Omega$ 上的 $C^1$ 函数, $C^\infty$ 函数, 或在 $\Omega$ 内有紧支集的 $C^\infty$ 函数组成. 尽管我们可能将这些算子中的任何一个都称为梯度, 它们却并不相同, 因而始终明确其定义域十分重要. 若说两个无界算子 $S$ 与 $T$ 相等, 则既是指 $D(S)=D(T)$, 也是指两算子的作用相同.}

% block: unbounded-sum-composition page:32 type:ordinary-paragraph
\BT{Two unbounded operators $S$ and $T$ from $X$ to $Y$ can be added to give an operator $S+T$, with domain $D(S+T)=D(S)\cap D(T)$. If $S$ is an unbounded operator $X\to Y$ and $T$ an unbounded operator $Y\to Z$, then we can similarly define the composition $T\circ S$ with domain $\{\,v\in D(S)\mid Sv\in D(T)\,\}$. One must be careful, since the sum or composition of densely defined operators need not be densely defined.}{从 $X$ 到 $Y$ 的两个无界算子 $S$ 和 $T$ 可以相加, 得到定义域为 $D(S+T)=D(S)\cap D(T)$ 的算子 $S+T$. 若 $S$ 是无界算子 $X\to Y$, $T$ 是无界算子 $Y\to Z$, 则类似地可以定义复合 $T\circ S$, 其定义域为 $\{\,v\in D(S)\mid Sv\in D(T)\,\}$. 必须注意, 稠密定义算子的和或复合不一定稠密定义.}

% block: null-range-graph-intro page:32 type:ordinary-paragraph
\BT{The \emph{null space} of an unbounded operator $T$, the \emph{range} of $T$, and the \emph{graph} of $T$ are each defined in the obvious way:}{无界算子 $T$ 的\emph{零空间}, $T$ 的\emph{像空间}和 $T$ 的\emph{图}分别按显然的方式定义:}

% block: null-range-graph-definitions page:32 type:equation
\[
\begin{aligned}
\mathcal N(T)&=\{\,x\in D(T)\mid Tx=0\,\}\subset X,\\
\mathcal R(T)&=\{\,Tx\mid x\in D(T)\,\}\subset Y,\\
\Gamma(T)&=\{\,(x,Tx)\mid x\in D(T)\,\}\subset X\times Y.
\end{aligned}
\]

% block: graph-norm-intro page:32 type:ordinary-paragraph
\BT{The function $v\mapsto(v,Tv)$ is a one-to-one mapping from the domain of $T$ onto its graph. Applying the norm in $X\times Y$ to $(v,Tv)$ defines a norm on $D(T)$. This norm, called the \emph{graph norm} on $D(T)$, is given by}{函数 $v\mapsto(v,Tv)$ 是从 $T$ 的定义域到其图的一一映射. 将 $X\times Y$ 中的范数应用于 $(v,Tv)$, 就在 $D(T)$ 上定义了一个范数. 这个范数称为 $D(T)$ 上的\emph{图范数}, 由下式给出:}

% block: graph-norm-equation page:32 type:equation
\[
\|v\|_{D(T)}^2:=\|v\|_X^2+\|Tv\|_Y^2,\qquad v\in D(T),
\]

% block: graph-inner-product-intro page:32 type:ordinary-paragraph
\BT{and is associated to the \emph{graph inner product}}{它与下述\emph{图内积}相联系:}

% block: graph-inner-product-equation page:32 type:equation
\[
\langle v,w\rangle_{D(T)}:=\langle v,w\rangle_X+\langle Tv,Tw\rangle_Y,\qquad v,w\in D(T).
\]

% block: graph-norm-strength page:32 type:ordinary-paragraph
\BT{Obviously the graph norm is stronger than the restriction to $T(D)$ of the $X$ norm, and commonly they are not equivalent. For example, in the example given above of the gradient operator, the graph norm is the usual $H^1$ norm, which is strictly stronger than the $L^2$ norm. We always have}{显然, 图范数强于 $X$ 范数在 $T(D)$ 上的限制, 而且通常两者不等价. 例如, 在前面的梯度算子例子中, 图范数就是通常的 $H^1$ 范数, 它严格强于 $L^2$ 范数. 总有}

% block: graph-boundedness-equation page:32 type:equation
\[
\|Tv\|_Y\leq\|v\|_{D(T)},\qquad v\in D(T),
\]

% block: graph-norm-completeness page:32 type:ordinary-paragraph
\BT{so the operator $T$ is a bounded operator from $D(T)$ to $Y$ when the former is endowed with the graph norm. The graph norm makes $D(T)$ an inner product space, or pre-Hilbert space, but it may not be complete.}{因此, 当前者配备图范数时, 算子 $T$ 是从 $D(T)$ 到 $Y$ 的有界算子. 图范数使 $D(T)$ 成为内积空间, 或预 Hilbert 空间, 但它未必完备.}%
% block: graph-norm-footnote page:32 type:footnote
\footnote{\BT{We could have instead defined the graph norm by $\|v\|_{D(T)}=\|v\|_X+\|Tv\|_Y$ without the squares. This definition gives an equivalent norm, but it would no longer be associated to an inner product.}{也可以不取平方, 将图范数定义为 $\|v\|_{D(T)}=\|v\|_X+\|Tv\|_Y$. 这个定义给出一个等价范数, 但它将不再与内积相联系.}}
